\documentclass[12pt]{article}
\usepackage[utf8]{inputenc}
\usepackage[T1]{fontenc}
\usepackage{hyperref}
\title{Celtiberian Language Resources for K-12 Curriculum Development}
\author{Compiled by Claude Sonnet 5.5 (claude-sonnet-5-5) for Andr\'as Kornai}
\date{}
\begin{document}
\maketitle

\section{Introduction}
Celtiberian (ISO 639-3 xce; Glottocode celt1247) is an extinct Celtic language of the Iberian Peninsula, attested in inscriptions written in the Celtiberian (Palaeohispanic) script and in some Latin-script texts. There are no native speakers and no descendant community. Botorrita I, III and IV are in Celtiberian script and make up most of the corpus, while Botorrita II is in Latin script. Sources differ on the dating of the Botorrita plaques (late 2nd or early 1st century BC).

\section{Primary Educational Resources}
I found no school curricula, primers or textbooks for Celtiberian.

\section{Linguistic and Reference Materials}
\begin{itemize}
\item Celtiberian language (Wikipedia): \url{https://en.wikipedia.org/wiki/Celtiberian_language}
\item Botorrita plaque (Wikipedia): \url{https://en.wikipedia.org/wiki/Botorrita_plaque}
\item Celtiberian script (Wikipedia): \url{https://en.wikipedia.org/wiki/Celtiberian_script}
\item David Stifter, Old Celtic Languages course notes on Celtiberian (University of Copenhagen): \url{https://rootsofeurope.ku.dk/kalender/arkiv_2012/celtic_spring/Keltiberisch_3_WS_2010.pdf}
\item \emph{Palaeohispanic Languages and Epigraphies} (Oxford University Press chapter): \url{https://academic.oup.com/book/36477/chapter/321084103}
\item Wolfgang Meid, \emph{Celtiberian Inscriptions} (Archaeolingua, 1994), catalogue record: \url{https://zenon.dainst.org/Record/000054064}
\end{itemize}

\section{Religious and Community Resources}
None found. There is no living community or religious use.

\section{Digital and Audio Resources}
No audio or software was found. The Lexvo entry for the language code is at \url{https://finto.fi/lexvo/en/page/xce}.

\section{Curriculum Development Context}
The corpus is small and only partly understood, and no one can speak the language. A K-12 heritage unit could use the Botorrita inscriptions and the Celtiberian script for history, archaeology and writing-systems lessons. Immersion use of the language is not possible, and any reconstructed phrases would be scholarly conjecture.
\end{document}